% Related work. Every reference was checked by web search; see paper/refs.bib for the two partly unverified entries.
\section{Related work}
\label{sec:related}

\paragraph{QAOA and its landscape} QAOA alternates a cost and a mixing layer, starting from the uniform superposition~\cite{farhi2014qaoa}. Its behaviour, parameter heuristics and non-adiabatic mechanism were analysed by Zhou et al.~\cite{zhou2020qaoa}, and a recent review covers its many variants~\cite{blekos2024review}. Barren plateaus make randomly initialised variational circuits hard to train~\cite{mcclean2018barren}. That is one motivation for informed initialisation, of both the parameters and the state.

\paragraph{Parameter initialisation and transfer} For fixed parameters, the QAOA objective concentrates over typical instances, so parameters transfer between instances~\cite{brandao2018concentration}. The transferability of optimal parameters between random graphs is explained by shared local subgraphs~\cite{galda2021transfer}. Annealing-inspired schedules initialise the parameters directly~\cite{sack2021tqa}, and graph neural networks have been trained to predict them~\cite{jain2022gnn}. These works learn the \emph{parameters}. We learn the \emph{initial state} and keep the parameter initialisation fixed and identical across arms.

\paragraph{Warm-starting} Egger et al.~\cite{egger2021warm} warm-start QAOA from a classical relaxation of each instance (a continuous relaxation for portfolio optimisation, Goemans--Williamson rounding for MaxCut). They introduce the regularised initial state and the matching mixer that we use, and show that warm-starting helps most at low depth. Tate et al.~\cite{tate2023sdp} warm-start from low-rank SDP solutions for MaxCut, with an advantage that shrinks as depth grows. Cain et al.~\cite{cain2022stuck} show that QAOA started from a single good classical string with an unmodified mixer makes little or no progress. Our $\varepsilon=0$ result, where a wrongly pinned bit can never flip, is a related failure, made exact by the warm mixer.

Prior work establishes that a given relaxation helps. We \emph{compare sources} of warm start under one protocol, measure what QAOA actually returns (exact best-of-$K$ ground-state sampling and operating cost), and identify the property of the initial state that determines it.

\paragraph{QUBO formulations and unit commitment} Lucas collected Ising formulations of many NP problems~\cite{lucas2014ising}. Classical UC is solved routinely by mixed-integer linear programming, with compact formulations that enforce minimum up and down times~\cite{carrion2006uc}. Machine learning has been used to predict good initial solutions and redundant constraints for large security-constrained UC instances~\cite{xavier2021learning}, which is the classical counterpart of a learned warm start.

\paragraph{Quantum methods for power systems} Reviews survey quantum computing for energy systems and power system optimisation~\cite{ajagekar2019energy,morstyn2024joule}, and for unit commitment specifically~\cite{hasanzadeh2026survey}. Koretsky et al.~\cite{koretsky2021qaoauc} adapted QAOA to UC with a classical minimiser for the continuous variables, and found classical solvers effective for their instances. Nikmehr et al.~\cite{nikmehr2022quc} decomposed UC across quantum and classical resources. We do not propose a new UC solver: UC serves as a structured, realistic testbed for what makes a warm start good.

\paragraph{Learning to optimise} Graph convolutional networks~\cite{kipf2017gcn} have been used inside exact combinatorial solvers~\cite{gasse2019exact}. Broader surveys cover machine learning for combinatorial optimisation~\cite{bengio2021tour}, learning to optimise~\cite{chen2022l2o} and amortised optimisation~\cite{amos2023amortized}, in which a model trained on a distribution of instances predicts solutions for new ones. Our learned warm start is an amortised predictor of this kind, and its out-of-distribution failure (Section~\ref{sec:res-neg}) is the usual price of amortisation.

\paragraph{QAOA against classical methods} Simulations comparing QAOA with classical MaxCut solvers place any crossover at hundreds of qubits or more~\cite{guerreschi2019maxcut}. Sampling-frequency analyses bound where an advantage could appear~\cite{lykov2023sampling}. Our compiled simulated-annealing baseline confirms, for our instances, that no advantage is in sight at \NLtwo--\NLfour\ qubits.
